\section{Discussion} \label{sec:discussion}

\subsection{"Lookahead" covariates help forecasting}

The experimental results consistently demonstrate that leveraging nearest neighbor information significantly improves the predictive accuracy of GBRT models. 
By identifying a historical subsequence of a nearest neighbor that is most similar to the current lookback window, the model gains access to a lookahead/template of what happened next in a similar past situation. 
This addresses a fundamental limitation of the standard sliding-window regression approach to forecasting, in which predictions are restricted to observations within the immediate lookback window (see Figure~\ref{fig:flattening}).
A simple model that incorporates only $NN_{\text{imm}}$ and $s$, namely \textsf{GBRT-NN-S}, always improves forecasting accuracy across all datasets, with few exceptions.
For \texttt{PeMSD7}, incorporating information about the nearest neighbor itself (i.e., $NN_{\text{last}}$, $NN_{\text{tar}}$, and $NN_{\text{cov}}$) improves the result.  

The included similarity $s$ can be treated as a confidence measure for the "Lookahead" covariates.
$NN_{\text{last}}$, $NN_{\text{tar}}$, and $NN_{\text{cov}}$ provides more information about the similar past situation.

\subsection{Raw values vs. other representations}

While Matrix Profile requires z-normalization for subsequence similarity computation, using the raw values of the immediate subsequence allows the model to recover lost magnitude information.

Pairwise differences are used for capturing trends.
However, the experiments show that this representation discards essential information required by GBRT to map predictors to responses effectively.

% \subsection{Date information dependency and history length}

\subsection{Include k-nearest neighbors information}

Expanding the feature set to include covariates derived from $k$-nearest neighbors ($k=2, 3$) can improve the performance for datasets, namely \texttt{Electricity} and \texttt{PeMSD7}. 
It can be seen from Table~\ref{tab:best_models} that models for $k$-nearest neighbors such as \textsf{GBRT-2NN\textsuperscript{++}-S}, \textsf{GBRT-3NN\textsuperscript{++}-S} and \textsf{GBRT-3NNC\textsuperscript{++}-S} appeared as the best models.
This suggests that the lookaheads from multiple historically similar occurrences provides a more robust predictor than a single neighbor, although it increases the dimensionality of the input vector.
To note, in \texttt{Electricity}, the best models are \textsf{GBRT} and \textsf{GBRT-NN} for RMSE in both scenarios, respectively.
It suggests that the above models can improve overall performance but cannot predict some outliers, resulting in large error terms $(y_i - \hat{y}_i)$, which RMSE is sensitive to.